\section{Síntesis y ampliación}

Lo que más vale la pena conservar de esta clase no es un speedup headline en particular, sino una forma de pensar que conecta architecture, statistics y systems: si de las estadísticas de los datos pueden obtenerse valores iniciales interpretables, es posible reducir el entrenamiento aleatorio; si la representación subyacente es fija, es posible almacenar en caché las comparaciones repetidas; si la tarea es muy estrecha y los datos se generan localmente, es posible trasladar el entrenamiento al edge. Cada paso trae nuevas restricciones: no se puede tomar solo el beneficio y descartar las premisas.

\subsection{Tres niveles de compresión: inicialización, reutilización del cómputo y ciclo de datos local}

El primer nivel es la initialization: bit-cipher proporciona una deterministic non-negative representation, y la generalized co-occurrence proporciona un warm start para la softmax layer. El segundo nivel es la computation reuse: los frozen embeddings permiten almacenar en caché las pairwise comparisons, y el dynamic context evita que las muestras cortas paguen el costo de la ventana más larga. El tercer nivel es el data loop: el edge controller recopila transcript/action pairs a partir de interacciones reales y cambia una tarea estrecha por la posibilidad de aprender con pocos datos.

\begin{importantbox}{La conclusión mínima confiable del curso}
\begin{enumerate}
  \item Para una softmax layer que cumple los supuestos, un data-derived warm start puede superar a un punto de partida aleatorio.
  \item El iterative training posterior al warm start sigue siendo valioso, sobre todo para la composition de varias capas y para la adaptación de la representación final.
  \item Las fixed representations permiten almacenar en caché parte de las comparaciones, pero la caché, la memoria y la invalidation deben contarse como costos del sistema.
  \item Un tiny local controller puede aprender acciones estrechas a partir de las interacciones con la aplicación, pero de ello no se puede deducir una capacidad de diálogo general.
\end{enumerate}
\end{importantbox}

\subsection{Checklist para reproducir los experimentos}

Al reproducir este tipo de trabajos, lo que más fácilmente se omite son los supuestos y los costos adicionales. La siguiente checklist debe publicarse junto con las gráficas de resultados.

\begin{lstlisting}[caption={Lista de reproducción y publicación de SAFFU/PLM}]
record_official_architecture_and_parameter_shapes()
verify_nonnegative_features_and_softmax_activation()
publish_explicit_initialization_formula_and_epsilon()
report_priming_number_distribution_and_sensitivity()
separate_stats_cache_backprop_and_evaluation_time()
compare_tuned_cold_warm_frozen_and_thawed_baselines()
document_cache_size_location_hit_rate_and_invalidation()
evaluate_heldout_users_paraphrases_noise_and_class_balance()
report_end_to_end_edge_latency_not_model_latency_only()
publish_code_data_splits_seeds_and_failure_cases()
\end{lstlisting}

\begin{warningbox}{Al ver «sin pretraining/backprop», pregunte de inmediato}
¿Qué capa es la que no lo necesita? ¿Cuántos datos usaron las estadísticas explícitas? ¿Se hace gradient tuning después? ¿El embedding está congelado? ¿La prueba mide training fit o held-out generalization? ¿Las unsupported layers se inicializan aleatoriamente? Si estas preguntas no tienen respuesta, el headline es mucho mayor que la evidencia.
\end{warningbox}

\subsection{Preguntas de autoevaluación}

\begin{multicols}{2}
\small
\begin{enumerate}
  \item ¿Qué relación hay, en cuanto a parametrización, entre la attention estándar $QK^\top$ y $XWX^\top$?
  \item ¿Por qué SAFFU sigue siendo una attention con parámetros y no una similitud sin parámetros?
  \item ¿Near-shallow describe la depth, o la actualización en línea y la dependencia de cómputo?
  \item ¿Por qué bit-cipher requiere $b\ge\lceil\log_2(V+1)\rceil$?
  \item ¿En qué se diferencia la generalized co-occurrence $H^\top Y$ de la coocurrencia tradicional en ventanas de palabras?
  \item ¿Por qué la softmax solution explícita requiere non-negative features?
  \item ¿Qué corrige el priming number $K$ en la fórmula?
  \item ¿Por qué los hidden targets de la attention son más difíciles de obtener que los decoder targets?
  \item ¿Por qué la gráfica de PPL con warm start no demuestra directamente que la attention de SAFFU sea mejor?
  \item ¿Qué pueden demostrar los resultados en MNIST y qué no pueden demostrar?
  \item ¿Por qué un frozen embedding afecta a la vez la velocidad de entrenamiento y la cache correctness?
  \item ¿Qué tipo de desperdicio reduce el packing y cuál el dynamic context?
  \item ¿Por qué la precisión de PLM no es FP16/BF16?
  \item ¿Qué optimiza cada una de las tres etapas warm--freeze--thaw?
  \item ¿Qué componentes incluye la latency de extremo a extremo de Le Potato?
  \item ¿Por qué la binary switch accuracy puede ser engañosa debido al class imbalance?
  \item ¿Qué conclusiones seguían siendo solo future work al terminar la clase?
\end{enumerate}
\end{multicols}

\subsection{Lecturas adicionales}

\begingroup
\small
\setlength{\itemsep}{2pt}
\begin{itemize}
  \item Williams and Zhao, \href{https://arxiv.org/abs/2311.07510}{\emph{Explicit Foundation Model Optimization with Self-Attentive Feed-Forward Neural Units}}: artículo principal sobre la arquitectura SAFFU, la inicialización explícita y la ablation.
  \item Williams and Zhao, \href{https://arxiv.org/abs/2311.07498}{\emph{Reducing the Need for Backpropagation and Discovering Better Optima With Explicit Optimizations of Neural Networks}}: solución explícita de softmax de una sola capa, MNIST y local multi-layer warm start.
  \item Zhao and Williams, \href{https://arxiv.org/abs/2311.11012}{\emph{Bit Cipher}}: deterministic bit vectors, co-occurrence aggregation y transfer experiments.
  \item Williams and Heidenreich, \href{https://arxiv.org/abs/2205.00148}{\emph{To Know by the Company Words Keep and What Else Lies in the Vicinity}}: Word2Vec softmax, co-occurrence factorization y la motivación del radial context.
  \item Frankle and Carbin, \href{https://arxiv.org/abs/1803.03635}{\emph{The Lottery Ticket Hypothesis}}: antecedente original para entender la motivación de esta clase de «reducir la lotería aleatoria»; no debe tomarse como conclusión experimental de SAFFU.
  \item Warstadt et al., \href{https://arxiv.org/abs/2310.00915}{\emph{Findings of the BabyLM Challenge}}: para revisar los requisitos concretos de developmentally plausible data scale y de la evaluation estándar.
\end{itemize}
\endgroup

\end{document}
